\begin{figure}[htbp]
			\centering
			\begin{tikzpicture}[
				scale=0.95,
				>=latex,
				trave/.style={line width=2pt},
				carico/.style={black, line width=0.8pt},
				forza/.style={black, line width=1pt},
				diagN/.style={black, line width=0.8pt},
				hatch/.style={gray!70, line width=0.4pt},
				quota/.style={black, line width=0.4pt},
				freccia/.style={black, line width=0.4pt, {Latex[length=1.5mm]}-{Latex[length=1.5mm]}}
				]
				
				% ── Parametri dei vincoli (convenzione C06/C07) ──
				\def\tw{0.15}    \def\thR{0.28}   \def\rn{0.06}
				\def\rw{0.05}    \def\bw{0.09}    \def\gw{0.22}
				\def\hh{0.10}
				
				% ── Geometria ──
				\def\Lell{4}
				\def\aa{2}
				\def\yq{-0.80}              % y della quota nel pannello (a)
				\def\ylab{\yq-0.7}          % y dei label (a) e (b) — allineati per costruzione
				
				% =========================================================
				% PANNELLO (a) — Schema strutturale
				% =========================================================
				\begin{scope}[xshift=0cm]
					
					% Trave
					\draw[trave] (0,0) -- (\Lell,0);
					
					% Cerniera in A
					\draw[thick] (0,0) -- (-\tw,-\thR) -- (\tw,-\thR) -- cycle;
					\draw[thick] (-\gw,-\thR) -- (\gw,-\thR);
					\foreach \x in {-0.16,-0.08,0,0.08,0.16}
					\draw[thin] (\x,-\thR) -- ++(-0.07,-\hh);
					
					% Carrello trasversale in B (asse di scorrimento orizzontale)
					\draw[thick] (\Lell,0) -- (\Lell-\tw,-\thR) -- (\Lell+\tw,-\thR) -- cycle;
					\draw[thick] (\Lell-\bw,-\thR-\rw-0.01) circle (\rw);
					\draw[thick] (\Lell+\bw,-\thR-\rw-0.01) circle (\rw);
					\draw[thick] (\Lell-\gw,-\thR-2*\rw-0.04) -- (\Lell+\gw,-\thR-2*\rw-0.04);
					\foreach \x in {-0.16,-0.08,0,0.08,0.16}
					\draw[thin] (\Lell+\x,-\thR-2*\rw-0.04) -- ++(-0.07,-\hh);
					
					% Anellini
					\draw[thick, fill=white] (0,0) circle (\rn);
					\draw[thick, fill=white] (\Lell,0) circle (\rn);
					
					% Etichette A, B, C
					\node[above left=-2pt,  font=\small] at (0,0)      {$A$};
					\node[above right=-2pt, font=\small] at (\Lell,0)  {$B$};
					\node[below,            font=\small] at (\aa,0.05) {$C$};
					
					% Carico distribuito q nel tratto AC
					
					% parametri geometrici del riquadro
					\def\hbase{0.22}    \def\hcar{0.50}
					
					% (1) frecce longitudinali sulla trave, in linea con la base del riquadro
					\pgfmathsetmacro{\yfre}{\hbase-0.04}
					\foreach \t in {0.10, 0.40, 0.70} {
						\pgfmathsetmacro{\px}{\aa*\t}
						\draw[->, carico] (\px, \yfre) -- (\px + 0.40, \yfre);
					}
					
					% (2) rettangolo soprastante con campitura verticale
					\foreach \x in {0.25, 0.50, ..., 1.75} {
						\draw[hatch] (\x, \hbase) -- (\x, \hbase+\hcar);
					}
					
					\draw[carico, line width=0.6pt] (0, \hbase) rectangle (\aa, \hbase+\hcar);
					
					% (3) etichetta del carico
					\node[right, font=\small] at (\aa+0.08, \hbase+0.5*\hcar) {$\bar{q}$};
					
					% Forza concentrata in B (orizzontale, verso destra)
					\def\Lf{0.85}
					\draw[->, forza] (\Lell, 0.0) -- (\Lell+\Lf, 0.0);
					\node[above, font=\small] at (\Lell+0.5*\Lf+0.1, -0.6) {$\bar{q}a$};
					
					% Quota a sotto la trave
					\def\yq{-0.80}
					\draw[quota] (0,\yq) -- (\aa,\yq);
					\draw[quota] (0,\yq+0.07) -- (0,\yq-0.07);
%					\draw[quota] (\aa,\yq+0.07) -- (\aa,\yq-0.07);
%					\node[above, font=\small] at (0.5*\aa, \yq-0.04) {$a$};
					
					% trattini di estremità (linee di riferimento)
					\draw[quota] (0,\yq+0.07) -- (0,\yq-0.07);
					\draw[quota] (\aa,\yq+0.07) -- (\aa,\yq-0.07);
					
					% linea di quota con frecce alle due estremità
					\draw[freccia] (0,\yq) -- (\aa,\yq);
					
					% etichetta
					\node[above, font=\small] at (0.5*\aa, \yq+0.02) {$a$};
					
					% Etichetta pannello
					\node[font=\small] at (0.75*\Lell, \yq-0.2) {(a)};
					
				\end{scope}
				
				% =========================================================
				% PANNELLO (b) — Diagramma di N(z)
				% =========================================================
				\begin{scope}[xshift=6cm]
					
					% Trave (clone)
					\draw[trave] (0,0) -- (\Lell,0);
					
					% Vincoli (clone esatto dal pannello a)
					\draw[thick] (0,0) -- (-\tw,-\thR) -- (\tw,-\thR) -- cycle;
					\draw[thick] (-\gw,-\thR) -- (\gw,-\thR);
					\foreach \x in {-0.16,-0.08,0,0.08,0.16}
					\draw[thin] (\x,-\thR) -- ++(-0.07,-\hh);
					\draw[thick] (\Lell,0) -- (\Lell-\tw,-\thR) -- (\Lell+\tw,-\thR) -- cycle;
					\draw[thick] (\Lell-\bw,-\thR-\rw-0.01) circle (\rw);
					\draw[thick] (\Lell+\bw,-\thR-\rw-0.01) circle (\rw);
					\draw[thick] (\Lell-\gw,-\thR-2*\rw-0.04) -- (\Lell+\gw,-\thR-2*\rw-0.04);
					\foreach \x in {-0.16,-0.08,0,0.08,0.16}
					\draw[thin] (\Lell+\x,-\thR-2*\rw-0.04) -- ++(-0.07,-\hh);
					\draw[thick, fill=white] (0,0) circle (\rn);
					\draw[thick, fill=white] (\Lell,0) circle (\rn);
					\node[left,  font=\small] at (0,-0.04)     {$A$};
					\node[right, font=\small] at (\Lell,-0.04) {$B$};
					
					% Diagramma N(z): scala 1 qa = 0.55 unità di figura
					% (convenzione: N positivo = trazione, disegnato sopra l'asse)
					\def\sN{0.55}
					\pgfmathsetmacro{\Nuno}{2*\sN}     % valore 2 qa
					\pgfmathsetmacro{\Ndue}{\sN}       % valore qa
					
					
					% dopo — tratto AC
					\foreach \x in {0.25, 0.50, ..., 1.75} {
						\pgfmathsetmacro{\h}{\Nuno - (\Nuno-\Ndue)*\x/\aa}
						\draw[hatch] (\x, 0) -- (\x, \h);
					}
					% dopo — tratto CB
					\foreach \x in {2.25, 2.50, ..., 3.75} {
						\draw[hatch] (\x, 0) -- (\x, \Ndue);
					}
					
					
					
					% Bordo del diagramma (linea spessa nera)
					\draw[diagN]
					(0, 0) -- (0, \Nuno) -- (\aa, \Ndue) -- (\Lell, \Ndue) -- (\Lell, 0);
					
					% Segno "+" (trazione) al centro dell'area
					\node[circle, draw, fill=white, thick, minimum size=10pt, inner sep=0pt, font=\large]	at (0.5*\Lell, 0.55*\sN-0.03) {$+$};
					
					% Etichette dei valori notevoli
					\node[below left,  font=\small] at (0,   \Nuno) {$2\bar{q}a$};
					\node[below right, font=\small] at (\Lell, \Ndue+0.2) {$\bar{q}a$};
					
					% Etichetta pannello
					%\node[font=\small\bfseries] at (0.5*\Lell, -1.80) {(b)};
					\node[font=\small] at (0.75*\Lell, \ylab+0.6) {(b)};
					
				\end{scope}
				
			\end{tikzpicture}
			\caption{Trave con carico longitudinale distribuito $\bar{q}$ nel tratto
				$AC$ e forza concentrata $\bar{q}a$ in $B$:
				(a)~schema strutturale;
				(b)~diagramma dello sforzo normale $N(z)$, in trazione su
				tutta la lunghezza.}
			\label{fig:stat_long_iso_01}
		\end{figure}
